\section{Paths of principal minors}\label{sec:minor-paths}

We prove the potential moment bound by following paths of nonsingular
principal minors. A tail estimate first controls the layers of large
height. On a short terminal window, either the potential is small at
some node, or a definite decrease of a second statistic forces many
steps that each have small conditional probability.

For a nonsingular matrix \(B\) and an integer \(s\ge0\), define
\[
W_s(B)=\sum_{p,e}(k_B(p,e)-s)_+\nu_p,
 \qquad x_+=\max\{x,0\}.
\]
In the layer diagram of Figure~\ref{fig:layers}, the factor
\((k_B(p,e)-s)_+\) counts the boxes in row \(e\) beyond the first
\(s\) columns; the sum weights them by \(\nu_p\).
When a tail occurs together with a nonsingularity requirement, set
\(W_s(B)=0\) on singular matrices.
The order bound in Equation~\eqref{eq:group-order} gives
\(0\le W_s(B)\le1/2\). For \(s\ge1\), we also have
\begin{equation}\label{eq:path-tail-difference}
 W_{s-1}(B)-W_s(B)=\sum_{\substack{p,e\\ k_B(p,e)\ge s}}\nu_p.
\end{equation}

\begin{lemma}[Tail recursion]\label{lem:tail-recursion}
Let \(C\to B\) be a transition of width \(h\in\{1,2\}\) between
nonsingular matrices. In the \(i\)-th successive quotient of \(G_C\),
write \(k_{i-1}(p,e)\) for the height before that quotient and
\(\mathcal R_i\) for its retained layers. For every integer \(s\ge h\),
\begin{equation}\label{eq:path-tail-recursion}
 W_s(B)\le W_{s+h}(C)
   +2\sum_{i=1}^{h}
       \sum_{\substack{(p,e)\in\mathcal R_i\\
                       k_{i-1}(p,e)\ge s+1-h}}\nu_p.
\end{equation}
\end{lemma}

\begin{proof}
Fix a prime \(p\), and put
\(F_D(p,t)=\sum_{e=1}^{t}k_D(p,e)\) for \(D=C,B\).
Take \(t\) to be the largest \(e\) for which \(k_B(p,e)>s\), or take
\(t=0\) when there is no such \(e\). Then
\[
 \sum_{e\ge1}(k_B(p,e)-s)_+=F_B(p,t)-st.
\]
For \(e\le t\), the height in \(G_B\) is at least \(s+1\).
The final successive quotient of \(G_C\) is also obtained from \(G_B\)
by \(h\) one-element quotients. The interlacing of heights under each
such quotient therefore gives
\[
 k_{i-1}(p,e)\ge k_B(p,e)-h\ge s+1-h\ge1
 \qquad (e\le t,\ 1\le i\le h).
\]
Let \(u_i(p,t)\) be the total drop in the first \(t\) layers at the
\(i\)-th quotient and put \(U_p(t)=\sum_i u_i(p,t)\). Each layer in
this range has positive starting height, and each of its drops is zero
or one. Thus
\[
 ht-U_p(t)=\sum_{i=1}^{h}
       \#\{e\le t:(p,e)\in\mathcal R_i\}.
\]
All the layers on the right have the height lower bound just displayed.
Lemma~\ref{lem:layer-inequality} now yields
\[
 \begin{split}
 F_B(p,t)-st
 &\le F_C(p,t)-(s+h)t+2\bigl(ht-U_p(t)\bigr)\\
 &\le \sum_{e\ge1}(k_C(p,e)-s-h)_+
       +2\sum_{i=1}^{h}
          \#\{e:(p,e)\in\mathcal R_i,\ k_{i-1}(p,e)\ge s+1-h\}.
 \end{split}
\]
Multiplying by \(\nu_p\) and summing over \(p\) proves the claim.
\end{proof}

Recall the scales \(R=\lceil\ell^2\rceil\) and \(\sigma=R^{-6}\) used
in Proposition~\ref{prop:retention-cost}. The next estimate applies
without a typicality hypothesis.

\begin{lemma}[Cost of heavy retained layers]\label{lem:heavy-retention-cost}
Let \(\mathcal F\) denote exposed information which fixes a nonsingular
binary matrix \(C\) of size \(m\in\mathcal I_N\) and a quotient map
\(\pi:\Z^m\twoheadrightarrow H\) with \(C\Z^m\subseteq\ker\pi\).
Suppose that, conditionally on \(\mathcal F\), a tested column \(v\) is
uniform on \(\bits^m\). If \(\mathcal R\) is the set of layers
retained on quotienting \(H\) by \(\pi(v)\), put
\[
 Q=\sum_{\substack{(p,e)\in\mathcal R\\ k_H(p,e)\ge R-1}}
           k_H(p,e)\nu_p.
\]
Fix \(0<\eta<1\) and set \(q_0=\eta/(100\log N)\). Uniformly for
\(q_0\le q\le1/2\), and for all sufficiently large \(N\) depending
on \(\eta\),
\begin{equation}\label{eq:path-heavy-retention}
 \Pr\{Q\in[q,2q)\mid\mathcal F\}\le 2^{-3qN/5}.
\end{equation}
The exposed information may include the preceding column of a width-two
transition, provided the tested column remains conditionally uniform.
\end{lemma}

\begin{proof}
Conditioning on \(\mathcal F\) fixes the quotient map and \(H\); it leaves \(v\)
uniform on \(\bits^m\).
For a fixed mask \(E\) of layers required to be retained,
Lemma~\ref{lem:retention-lattice} supplies the necessary condition
\[
 v\in K_E,\qquad
 C\Z^m\subseteq K_E\subseteq\Z^m,\qquad
 [\Z^m:K_E]=N^{q(E)N},\qquad
 q(E)=\sum_{(p,e)\in E}k_H(p,e)\nu_p.
\]
We first show that this condition costs at least
\(2q(E)N/3\) bits of probability.

Write \(g_1,\ldots,g_m\) for the images of the coordinate vectors in
\(\Z^m/K_E\). Choose a maximal subset \(g_{i_1},\ldots,g_{i_j}\) that
is dissociated, meaning that all of its subset sums are distinct.
After the other bits of \(v\) are fixed, at most one choice of the
\(j\) selected bits can give any prescribed group element. Hence every
atom of the image of \(v\), in particular the zero atom, has
probability at most \(2^{-j}\).

Maximality also implies that every \(g_a\) is a combination of the
selected elements with coefficients in \(\{0,1,-1\}\): otherwise
adjoining \(g_a\) would preserve distinctness of subset sums. Choose
such representations, using the identity representations at the
selected indices, and arrange their coefficients as the columns of a
matrix \(M\in\Z^{j\times m}\). This matrix has rank \(j\). The selected
elements generate \(\Z^m/K_E\), so they define a surjection
\(\theta:\Z^j\to\Z^m/K_E\). Every column of \(MC\) belongs to
\(\ker\theta\), because the columns of \(C\) belong to \(K_E\).
Since \(C\) is nonsingular, \(MC\) still has rank \(j\). Choose \(j\)
independent columns of \(MC\). Each of their entries has absolute
value at most \(m\le N\). Their lattice is contained in
\(\ker\theta\), and Hadamard's inequality gives
\[
 |\Z^m/K_E|=[\Z^j:\ker\theta]
 \le (N\sqrt j)^j\le N^{3j/2}.
\]
For \(j=0\), the group is trivial and the same conclusion has its
usual empty-product meaning. It follows that
\begin{equation}\label{eq:path-fixed-mask}
 \Pr\{v\in K_E\mid\mathcal F\}\le2^{-j}\le2^{-2q(E)N/3}.
\end{equation}

We next bound the number of possible masks consisting of all retained
layers of height at least \(R-1\), when their value of \(Q\) lies in
\([q,2q)\). At a fixed prime the retained layers on each plateau form
a prefix. Let \(d_p\) be the number of plateaus of \(H\) at \(p\).
Their distinct positive heights imply
\[
 d_p^2\le 2\sum_e k_H(p,e).
\]
Moreover,
\(\sum_{p,e}k_H(p,e)\le\log_2|H|=O(N\log N)\).
There are at most \(N^{1/2}\) primes \(p\le N^{1/2}\), so
Cauchy--Schwarz bounds the number of their plateaus by
\[
 \sum_{p\le N^{1/2}}d_p
   =O\bigl(N^{3/4}(\log N)^{1/2}\bigr).
\]
Each prefix length has \(O(N\log N)\) possible values and costs
\(O(\log N)\) bits. The full cost for these primes is consequently
\[
 O\bigl(N^{3/4}(\log N)^{3/2}\bigr)
     =o(N/\log N)=o_\eta(qN).
\]
Here the final estimate is uniform for \(q\ge q_0\).

For \(p>N^{1/2}\), each eligible layer has, for large \(N\),
\[
 k_H(p,e)\nu_p\ge\frac{R-1}{2N}\ge\frac{R}{3N}.
\]
The total height times weight in \(H\) is at most \(1/2\), so there are
at most \(3N/(2R)\) eligible layers in this range. A mask with \(Q<2q\)
selects at most \(6qN/R\) of them. The usual binomial estimate, using
the bound \(2^a\) for all subsets when the permitted number is a
fixed fraction of the \(a\) available layers, gives a logarithmic
count of at most
\[
 O\left(\frac{qN}{R}\bigl(1+\log(1/q)\bigr)\right)=o_\eta(qN).
\]
Indeed \(qN/R\ge\eta N/(100R\log N)\to\infty\), so integer rounding
does not affect this estimate, and
\[
 1+\log(1/q)\le 1+\log(100\log N/\eta)=O_\eta(\ell),
 \qquad R\asymp\ell^2.
\]
Combining the two prime ranges, the mask entropy is \(o_\eta(qN)\)
uniformly in the indicated interval. For every mask being counted,
\(q(E)=Q\ge q\). The union bound with
Equation~\eqref{eq:path-fixed-mask} is therefore
\[
 \Pr\{Q\in[q,2q)\mid\mathcal F\}
 \le 2^{-(2/3-o_\eta(1))qN}\le2^{-3qN/5}.
\]
The conditioning argument uses only that \(\mathcal F\) fixes \(C,\pi\)
and leaves \(v\) uniform, so it also applies when \(\mathcal F\) includes
a preceding column.
\end{proof}

\begin{proof}[Proof of Proposition~\ref{prop:path-potential}]
Fix \(0<\epsilon<1/10\), and put
\[
 \epsilon'=\epsilon/20,\qquad
 J_*=2\lceil N^{3/5}\rceil,\qquad
 L_*=\lceil R^{12}\rceil.
\]
Here \(J_*\) is the length of the long path, and \(L_*\) is the length
of its terminal window. The following comparisons will be used throughout:
\begin{equation}\label{eq:path-scales}
 \begin{gathered}
 2J_*+1<N^{7/10},\qquad
 L_*\asymp\ell^{24}=o(N^{3/5}),\qquad
 J_*-L_*>\lceil N^{3/5}\rceil,\\
 L_*\ell^{-60}=O(\ell^{-36})=o(1),\qquad
 \frac{L_*\sigma}{R^3}\asymp R^3\longrightarrow\infty,\qquad
 \frac{N\sigma}{\log R}=\frac{N}{R^6\log R}\longrightarrow\infty.
 \end{gathered}
\end{equation}
They hold for all sufficiently large \(N\), since \(R\asymp\ell^2\)
and every fixed power of \(\ell\) is \(o(N^a)\) for every fixed \(a>0\).
The first line keeps every node in \(\mathcal I_N\) and leaves enough
preceding steps to drive the tail threshold past \(N^{3/5}\). The second
line keeps the potential-step error small over the terminal window, makes
the accumulated costly-step exponent diverge, and makes that cost dominate
the entropy of choosing the costly positions.

For \(m\in\mathcal I_N\), let \(\mathcal T_m\) be the set of
nonsingular matrices satisfying the estimates in
Propositions~\ref{prop:prime-corank} and \ref{prop:retention-cost},
and those in Propositions~\ref{prop:admissible-potential} and
\ref{prop:potential-step} at precision \(A=60\), including both
widths in the latter proposition. These finitely many properties can
be chosen so that, for some \(f(N)\to\infty\),
\[
 \Pr\{\det\mathbf B_m\ne0,\ \mathbf B_m\notin\mathcal T_m\}
 \le2^{-Nf(N)}
 \qquad(m\in\mathcal I_N).
\]
Write \(m_*=N-1\). Consider all principal submatrices of
\(\mathbf B_{m_*}\) obtained by deleting at most \(2J_*\) indices.
Their sizes lie in \(\mathcal I_N\) by
Equation~\eqref{eq:path-scales}, and their number is at most
\[
 \sum_{d\le2J_*}\binom{m_*}{d}
 \le2^{O(N^{3/5}\log N)}=2^{o(N)}.
\]
Let \(\mathcal E_N\) be the event that any nonsingular one of these
submatrices fails its typical properties. Each fixed submatrix is a
binary matrix, so the union bound gives
\begin{equation}\label{eq:path-typical-exception}
 \Pr(\mathcal E_N)
 \le 2^{-N(f(N)-O(N^{-2/5}\log N))}=2^{-\omega(N)}.
\end{equation}
Since \(2^{N\Psi}\le2^N\), its contribution to the moment is at most
\[
 2^N\Pr(\mathcal E_N)
 \le2^{-N(f(N)-1-O(N^{-2/5}\log N))}=2^{-\omega(N)}.
\]

Starting from a nonsingular realization on \([m_*]\), apply
Lemma~\ref{lem:gated-deletion} backwards for exactly \(J_*\) steps.
It gives a path of nonsingular principal minors, each deletion having
width one or two, and every deletion of width two satisfying the gate
when read as a forward transition. At most \(2J_*\) indices are
deleted, so all these operations stay in the prescribed size range.
On \(\mathcal E_N^c\), all of their nodes also belong to their
respective sets \(\mathcal T_m\).

For clarity, the path just obtained is chosen backwards as a function
of the realized terminal matrix. We use it only to cover the terminal
event by path events. Let \(\Pi_N\) be the collection of all
\emph{fixed} index paths
\[
 I_0\subset I_1\subset\cdots\subset I_{J_*}=[m_*],
 \qquad h_t=|I_t\setminus I_{t-1}|\in\{1,2\}.
\]
An ordering of the two indices in a pair is fixed once for each path.
At each backward step there are at most \(2N^2\) choices, whence
\[
 |\Pi_N|\le(2N^2)^{J_*}
       =2^{O(N^{3/5}\log N)}=2^{o(N)}.
\]
For a fixed \(\pi\in\Pi_N\), put \(B_t=\mathbf B_{m_*}[I_t]\). Let
\(\Gamma_t\) be the event \(\mathcal G_{h_t}\) from
Equation~\eqref{eq:extension-gate} for step \(t\). Define
\[
 \mathcal A_\pi
   =\bigcap_{t=0}^{J_*}\{B_t\in\mathcal T_{|I_t|}\}
      \ \cap\ \bigcap_{t=1}^{J_*}\Gamma_t.
\]
The backward construction shows that every nonsingular terminal
matrix outside \(\mathcal E_N\) lies in at least one
\(\mathcal A_\pi\). Consequently
\begin{equation}\label{eq:path-cover}
 \E\!\left[\ind_{\{\det\mathbf B_{m_*}\ne0\}}\,
                 2^{N\Psi(\mathbf B_{m_*})}\right]
 \le2^N\Pr(\mathcal E_N)
       +\sum_{\pi\in\Pi_N}
          \E\!\left[\ind_{\mathcal A_\pi}2^{N\Psi(B_{J_*})}\right].
\end{equation}
We now bound a summand for a fixed \(\pi\). Expose the entries of
\(B_0\) first and then expose each transition forwards. Conditional
on the entries already exposed, its cross columns are independent
uniform bit vectors on the starting indices, and its new symmetric
block consists of independent bits. Within a pair the second cross
column is still uniform after the first is exposed. We sum over fixed
paths before making any of these exposures; we never condition this
forward process on the backward choice of a path. The future
nonsingularity, typicality, and gate requirements are imposed only
through indicators.

Call the last \(L_*\) steps and their \(L_*+1\) nodes the terminal window.
We will prove the three weighted bounds in Table~\ref{tab:path-contributions}.
\begin{table}[ht]
\centering
\renewcommand{\arraystretch}{1.15}
\begin{tabular}{@{}p{0.52\textwidth}p{0.42\textwidth}@{}}
\toprule
Covered event & Terminal weighted contribution\\
\midrule
Some window node has \(W_R>\epsilon'\)
 & \(2^{N-\Omega_\epsilon(NR)}\)\\
Some window node has \(\Psi\le\epsilon\)
 & \((L_*+1)2^{N(\epsilon+L_*\ell^{-60})}\)\\
Every window node has \(W_R\le\epsilon'\) and \(\Psi>\epsilon\)
 & \(2^{N-\Omega_\epsilon(NR^3)}\)\\
\bottomrule
\end{tabular}
\caption{Bounds to prove for a cover of the terminal weighted contribution
of one fixed path. Each row includes the indicator of \(\mathcal A_\pi\),
the indicator of its displayed event, and the weight \(2^{N\Psi(B_{J_*})}\).
The first two events may overlap.}
\label{tab:path-contributions}
\end{table}

\paragraph{Large tails in the window.}
At each cross column of the
fixed path, form \(Q\) as in Lemma~\ref{lem:heavy-retention-cost}
using the quotient immediately before that column. Put
\[
 q_{\min}=\frac{\epsilon'}{100\log N},\qquad
 \widetilde Q=
 \begin{cases}
 Q,&Q\ge q_{\min},\\
 0,&Q<q_{\min}.
 \end{cases}
\]
For an unconditional exposure on which the starting matrix \(C\) of
a step is singular, define \(\widetilde Q=0\) for its cross columns.
This convention permits estimates on the entire probability space.
For a nonsingular start, \(Q\le1/2\) by the group order bound.
Apply Lemma~\ref{lem:heavy-retention-cost} with \(\eta=\epsilon'\)
to the dyadic values \(q_j=2^j q_{\min}\) up to \(1/2\). Conditional
on everything exposed before a cross column,
\[
 \begin{split}
 \E\!\left[2^{N\widetilde Q/4}\mid
                    \hbox{previous exposures}\right]
 &\le1+\sum_{j:q_j\le1/2}
          2^{Nq_j/2}\,2^{-3Nq_j/5}\\
 &\le1+\sum_{j:q_j\le1/2}2^{-Nq_j/10}
 \le1+\delta_N,\qquad
 \delta_N=2^{-\epsilon'N/(2000\log N)}.
 \end{split}
\]
The last inequality holds for large \(N\), since
\(Nq_{\min}\to\infty\); the dyadic sum is at most twice its first term
and is then bounded by the displayed \(\delta_N\). It holds also for
a singular start by the convention above. There are
\(d_\pi=\sum_t h_t\le2J_*=O(N^{3/5})\) cross columns. Successive
conditioning, including between the two columns of a pair, therefore
gives the multiplicative bound
\begin{equation}\label{eq:path-heavy-moment}
 \E\,2^{(N/4)\sum_{a=1}^{d_\pi}\widetilde Q_a}
 \le(1+\delta_N)^{d_\pi}
 \le\exp\!\bigl(O(N^{3/5}\delta_N)\bigr)=1+o(1).
\end{equation}
In particular the accumulated error in its logarithm is
\(O(N^{3/5}2^{-\epsilon'N/(2000\log N)})=o(1)\).

On \(\mathcal A_\pi\), suppose \(W_R(B_t)>\epsilon'\) at a node
\(J_*-L_*\le t\le J_*\). Starting at this node, iterate
Lemma~\ref{lem:tail-recursion} backwards, increasing \(s\) by the
width at every step. Stop as soon as \(s>N^{3/5}\). This takes at most
\(\lceil N^{3/5}\rceil\) steps, so the available preceding path is
long enough by Equation~\eqref{eq:path-scales}. Every invocation has
\(s\ge R\ge2\), as required. At the last node the tail is zero:
Proposition~\ref{prop:prime-corank} gives
\[
 k_B(p,e)\le k_B(p,1)=\corank_{\F_p}(B\bmod p)\le N^{3/5}
\]
at every typical nonsingular node and every prime \(p\).

Number the backwards steps from the chosen node by \(j=1,2,\ldots,r\).
Before step \(j\), the current value of \(s\) is \(R\) plus the widths
of its \(j-1\) predecessors. If the current width is \(h\le2\), each
charged layer in Equation~\eqref{eq:path-tail-recursion} thus has
height at least
\[
 s+1-h\ge R+j-2.
\]
The charge from one quotient at this step is bounded by
\(2Q/(R+j-2)\), because \(Q\) includes height times weight for all
retained layers of height at least \(R-1\). There are at most two
quotients per step. The terms with \(Q<q_{\min}\) contribute at most
\[
 4q_{\min}\sum_{j=1}^{r}\frac1{R+j-2}
 \le4q_{\min}(1+\log N)\le\epsilon'/2
\]
for large \(N\). Here \(r\le\lceil N^{3/5}\rceil\) and the last
comparison follows from the definition of \(q_{\min}\).
The remaining charges must therefore exceed \(\epsilon'/2\).
Since every denominator is at least \(R-1\), this implies
\[
 \sum_{a=1}^{d_\pi}\widetilde Q_a
       \ge\frac{\epsilon'(R-1)}4.
\]
The same condition on the sum over the whole path is necessary if
any node of the window has \(W_R>\epsilon'\). By Markov's inequality
and Equation~\eqref{eq:path-heavy-moment},
\begin{equation}\label{eq:path-large-tail}
 \begin{split}
 &\Pr\!\left(\mathcal A_\pi,\
       \max_{J_*-L_*\le t\le J_*}W_R(B_t)>\epsilon'\right)\\
 &\qquad\le
 2^{-\epsilon'N(R-1)/16}\,
      \exp\!\bigl(O(N^{3/5}\delta_N)\bigr)
 =2^{-\Omega_\epsilon(NR)}=2^{-\omega(N)}.
 \end{split}
\end{equation}
Multiplication by the terminal weight, which is at most \(2^N\),
still gives \(2^{-\omega(N)}\), because \(R\to\infty\).

\paragraph{A small potential in the window.}
Proposition~\ref{prop:potential-step} at precision \(60\)
states, for either width and for a starting matrix \(C\in\mathcal T_m\),
\[
 \E\!\left[\ind_{\{\det B\ne0\}}\ind_{\Gamma}\,
                   2^{N\Psi(B)}\mid C\right]
       \le2^{N(\Psi(C)+\ell^{-60})},
\]
where \(\Gamma\) is the gate for width two and the certain event for
width one. Adding a typicality indicator at the endpoint only
decreases the left side. Successive conditioning from a node \(t\)
to the terminal node consequently bounds the conditional weighted
expectation with all the intervening requirements by
\[
 2^{N\Psi(B_t)}\,2^{N(J_*-t)\ell^{-60}}
\]
whenever \(B_t\in\mathcal T_{|I_t|}\). This is an application to the
fixed forward exposure, with no conditioning on future requirements.
For each \(t\) in the window, on \(\Psi(B_t)\le\epsilon\) it is at
most \(2^{N(\epsilon+L_*\ell^{-60})}\). Dropping all requirements
before \(t\) and summing over its \(L_*+1\) possible values gives
\begin{equation}\label{eq:path-low-window}
 \begin{split}
 \E\!\left[\ind_{\mathcal A_\pi}
       \ind_{\{\min_{J_*-L_*\le t\le J_*}\Psi(B_t)\le\epsilon\}}
                  2^{N\Psi(B_{J_*})}\right]
 &\le(L_*+1)\,2^{N(\epsilon+L_*\ell^{-60})}\\
 &=2^{N(\epsilon+O(\ell^{-36})+O(\log R/N))}.
 \end{split}
\end{equation}
This records both the accumulated transition error and the cost of
choosing the node.

\paragraph{The remaining window event.}
Consider the event
\[
 \mathcal H_\pi=\mathcal A_\pi
  \cap\left\{\max_{J_*-L_*\le t\le J_*}W_R(B_t)\le\epsilon'\right\}
  \cap\left\{\min_{J_*-L_*\le t\le J_*}\Psi(B_t)>\epsilon\right\}.
\]
At any node \(B\) of this event, \(\phi(k)\le2k\) and
\(\phi(0)=\phi(1)=0\) show that
\(\sum_{k_B\ge2}k_B\nu_p\ge\epsilon/2\). On the other hand,
\[
 \sum_{\substack{p,e\\ k_B(p,e)>2R}}k_B(p,e)\nu_p
 \le2W_R(B)\le2\epsilon'.
\]
Dividing the contribution from \(2\le k_B\le2R\) by \(2R\) gives
\begin{equation}\label{eq:path-window-two}
 \sum_{\substack{p,e\\k_B(p,e)\ge2}}\nu_p
 \ge\frac{\epsilon/2-2\epsilon'}{2R}
 \ge\frac{\epsilon}{10R}.
\end{equation}
If in addition \(\Psi(B)>1-1/R\), the contribution to \(\Psi(B)\)
from height two is at most \(\beta/4=7/20\). Indeed,
\(2\sum_{k_B=2}\nu_p\le1/2\) by the order bound. Thus
\[
 2\sum_{\substack{p,e\\k_B(p,e)\ge3}}k_B(p,e)\nu_p
  =\Psi(B)-\beta\sum_{\substack{p,e\\k_B(p,e)=2}}\nu_p
  >1-\frac1R-\frac{\beta}{4}.
\]
For large \(N\) the height times weight on the left after division by
two is at least \(1/4\). Since \(\epsilon'<1/200\), deleting the
layers above \(2R\) and dividing by \(2R\) yields
\begin{equation}\label{eq:path-window-three}
 \sum_{\substack{p,e\\k_B(p,e)\ge3}}\nu_p
 \ge\frac{1/4-2\epsilon'}{2R}\ge\frac1{20R}.
\end{equation}

For each quotient at step \(t\), let \(V_{t,i}\) be the sum of the
weights \(\nu_p\) of all of its retained layers. On
\(\mathcal A_\pi\), call a step in the window costly if some
\(V_{t,i}\ge\sigma\), or if it has width two and its starting node
\(C=B_{t-1}\) satisfies \(\Psi(C)\le1-1/R\). To price this
description without conditioning on \(\mathcal A_\pi\), define an
exposure event \(D_t\) as follows: the starting node belongs to its
set \(\mathcal T_m\), and either some \(V_{t,i}\ge\sigma\), or the
step has width two, \(\Psi(C)\le1-1/R\), and its gate \(\Gamma_t\)
holds. Every costly step of a required path satisfies \(D_t\).

Let \(\mathcal F_{t-1}\) be the sigma-field generated by all entries
exposed through \(C\).
Proposition~\ref{prop:retention-cost}, applied to each column and
conditioned on a preceding column when needed, gives probability
at most \(2^{-N\sigma/2}\) for each event \(V_{t,i}\ge\sigma\).
For the other alternative, the two cross columns are conditionally
independent and the gate requires both to belong to \(\mathcal S(C)\).
Proposition~\ref{prop:admissible-potential} gives
\[
 \begin{split}
 \Pr(\Gamma_t\mid\mathcal F_{t-1})
 &=\left(\frac{|\mathcal S(C)|}{2^m}\right)^2\\
 &\le2^{-2N(1-\Psi(C)-\ell^{-60}-(1-m/N))}.
 \end{split}
\]
Here \(1-m/N\le N^{-3/10}\). If
\(\Psi(C)\le1-1/R\), then
\(\ell^{-60}+N^{-3/10}=o(R^{-1})\) and
\(\sigma=R^{-6}=o(R^{-1})\), so for large \(N\) this probability is
at most \(2^{-N/R}\le2^{-N\sigma/2}\). If the start is not in
\(\mathcal T_m\), the event \(D_t\) is empty. The union bound therefore
gives, uniformly in the revealed history,
\begin{equation}\label{eq:path-costly-step}
 \Pr(D_t\mid\mathcal F_{t-1})
 \le p_N,\qquad p_N=3\,2^{-N\sigma/2}.
\end{equation}
In particular, for any fixed set \(T\) of step positions, ordinary
successive conditioning in the forward exposure gives
\begin{equation}\label{eq:path-costly-set}
 \Pr\left(\bigcap_{t\in T}D_t\right)\le p_N^{|T|}.
\end{equation}
This estimate imposes no conditioning on the other path requirements.

\paragraph{Descent at noncostly steps.}
To force many costly steps on \(\mathcal H_\pi\), we combine two tails.
At width one, the smaller \(R^{-2}W_1\) term uses the mass of layers of
height at least two. At width two, \(W_2\) uses the stronger gap between
\(W_2\) and \(W_4\), while the small coefficient prevents the \(W_1\)
term from offsetting that gap. Define
\[
 Z(B)=W_2(B)+R^{-2}W_1(B).
\]
The order bound implies \(0\le Z(B)\le(1+R^{-2})/2\le1\).
At a noncostly step all its \(V_{t,i}\) are less than \(\sigma\).
For a step of width one, Lemma~\ref{lem:tail-recursion} gives
\[
 W_2(B)\le W_3(C)+2\sigma\le W_2(C)+2\sigma,
 \qquad W_1(B)\le W_2(C)+2\sigma.
\]
Equations~\eqref{eq:path-tail-difference} and
\eqref{eq:path-window-two} then show
\[
 Z(B)-Z(C)
 \le2\sigma(1+R^{-2})
       -R^{-2}\bigl(W_1(C)-W_2(C)\bigr)
 \le2\sigma(1+R^{-2})-\frac{\epsilon}{10R^3}
 \le-\frac{\epsilon}{20R^3}
\]
for large \(N\), since \(\sigma=R^{-6}=o_\epsilon(R^{-3})\).
At a noncostly step of width two, its starting potential is
\(\Psi(C)>1-1/R\). Lemma~\ref{lem:tail-recursion} gives
\[
 W_2(B)\le W_4(C)+4\sigma.
\]
Also
\[
 W_2(C)-W_4(C)
 =\sum_{\substack{p,e\\k_C(p,e)\ge3}}
        \min\{k_C(p,e)-2,2\}\nu_p
 \ge\frac1{20R}
\]
by Equation~\eqref{eq:path-window-three}. Using \(W_1(B)\le1/2\),
we obtain
\[
 Z(B)-Z(C)
 \le-\frac1{20R}+4\sigma+\frac1{2R^2}
 \le-\frac1{40R}\le-\frac{\epsilon}{20R^3}
\]
for all sufficiently large \(N\). The middle comparison uses
\(4R^{-6}+(2R^2)^{-1}=o(R^{-1})\).
Thus every noncostly step decreases \(Z\) by at least
\[
 \delta=\frac{\epsilon}{20R^3}.
\]

Let \(n_{\mathrm{cost}}\) be the number of costly steps in this window.
A costly step can increase \(Z\) by at most one. Telescoping the \(L_*\) changes and
using \(0\le Z\le1\) gives
\[
 -1\le Z(B_{J_*})-Z(B_{J_*-L_*})
       \le n_{\mathrm{cost}}-(L_*-n_{\mathrm{cost}})\delta,
 \qquad
 n_{\mathrm{cost}}\ge\frac{L_*\delta-1}{1+\delta}.
\]
Now \(\delta\to0\) and
\(L_*\delta\asymp_\epsilon R^9\to\infty\). Consequently, for large
\(N\),
\[
 n_{\mathrm{cost}}\ge \frac{L_*\epsilon}{100R^3}.
\]
Set \(t_*=\lceil L_*\epsilon/(100R^3)\rceil\). The event
\(\mathcal H_\pi\) therefore implies \(D_t\) at at least \(t_*\)
positions in the window. Equations~\eqref{eq:path-costly-step} and
\eqref{eq:path-costly-set} yield
\[
 \Pr(\mathcal H_\pi)\le\binom{L_*}{t_*}p_N^{t_*}.
\]
The entropy of choosing these positions satisfies
\[
 \log_2\binom{L_*}{t_*}
 \le t_*\log_2(eL_*/t_*)=O_\epsilon(t_*\log R).
\]
Together with the factor \(3^{t_*}\), this is
\(o(N\sigma t_*)\), because \(N\sigma/\log R\to\infty\).
In particular, for all sufficiently large \(N\),
\begin{equation}\label{eq:path-high-window}
 \Pr(\mathcal H_\pi)
 \le2^{-N\sigma t_*/3}
 \le2^{-\epsilon N L_*\sigma/(300R^3)}
 =2^{-\Omega_\epsilon(NR^3)}
 =2^{-\omega(N)}.
\end{equation}
This bound uses \(L_*\sigma/R^3\asymp R^3\to\infty\).
The corresponding weighted expectation is at most
\(2^{N-\Omega_\epsilon(NR^3)}=2^{-\omega(N)}\).

Equations~\eqref{eq:path-large-tail}, \eqref{eq:path-low-window}, and
\eqref{eq:path-high-window} prove the three weighted bounds in
Table~\ref{tab:path-contributions}. Its row events cover all outcomes in
\(\mathcal A_\pi\), so their sum gives, uniformly over all fixed index paths,
\[
 \E\!\left[\ind_{\mathcal A_\pi}2^{N\Psi(B_{J_*})}\right]
 \le2^{N(\epsilon+o_\epsilon(1))}.
\]
Finally use Equation~\eqref{eq:path-cover}. The index-path sum adds
\(O(N^{3/5}\log N)=o(N)\) to the logarithm, and
Equation~\eqref{eq:path-typical-exception} gives a negligible
exceptional contribution even after the weight \(2^N\). Hence
\[
 \E\!\left[\ind_{\{\det\mathbf B_{N-1}\ne0\}}\,
                 2^{N\Psi(\mathbf B_{N-1})}\right]
 \le2^{N(\epsilon+o_\epsilon(1))}.
\]
For each fixed \(\epsilon\) this estimate holds for all sufficiently
large \(N\). Letting \(\epsilon\) tend to zero after \(N\) proves
Proposition~\ref{prop:path-potential}.
\end{proof}

\begin{proof}[Proof of Theorem~\ref{thm:main}]
By the implication established in Section~\ref{sec:framework},
Propositions~\ref{prop:admissible-potential} and \ref{prop:path-potential}
give Equation~\eqref{eq:column-target}:
\(\E|\mathcal S(\mathbf B_{N-1})|\le2^{o(N)}\).

Lemmas~\ref{lem:corank-reduction} and \ref{lem:column-reduction} now
give
\[
 \Pr(\det\mathbf B_N=0)
 \le2^{o(N)}N2^{-N}\E|\mathcal S(\mathbf B_{N-1})|
       +2^{-\Omega(N^{6/5})}
 \le2^{-N+o(N)}.
\]
A specified row of \(\mathbf B_N\) is identically zero with probability
\(2^{-N}\), so the reverse bound follows from that event. Therefore
\(\Pr(\det\mathbf B_N=0)=2^{-N+o(N)}\).
Finally Lemma~\ref{lem:sign-bit}, with \(N=n-1\), yields
\[
 \Pr(\det A_n=0)=2^{-(n-1)+o(n)}
                =2^{-n+o(n)}
                =\left(\frac12+o(1)\right)^n.
\]
\end{proof}
